%document
\documentclass[10pt]{beamer}
%theme
\usetheme{metropolis}
% packages
\usepackage{color}
\usepackage{listings}
\usepackage[ngerman]{babel}
\usepackage[utf8]{inputenc}
\usepackage{multicol}

\usepackage{hyperref}
\hypersetup{colorlinks=true, linkcolor=blue, urlcolor=red}

% color definitions
\definecolor{mygreen}{rgb}{0,0.6,0}
\definecolor{mygray}{rgb}{0.5,0.5,0.5}
\definecolor{mymauve}{rgb}{0.58,0,0.82}

\lstset{
    backgroundcolor=\color{white},
    % choose the background color;
    % you must add \usepackage{color} or \usepackage{xcolor}
    basicstyle=\footnotesize\ttfamily,
    % the size of the fonts that are used for the code
    breakatwhitespace=false,
    % sets if automatic breaks should only happen at whitespace
    breaklines=true,                 % sets automatic line breaking
    captionpos=b,                    % sets the caption-position to bottom
    commentstyle=\color{mygreen},    % comment style
    % deletekeywords={...},
    % if you want to delete keywords from the given language
    extendedchars=true,
    % lets you use non-ASCII characters;
    % for 8-bits encodings only, does not work with UTF-8
    frame=single,                    % adds a frame around the code
    keepspaces=true,
    % keeps spaces in text,
    % useful for keeping indentation of code
    % (possibly needs columns=flexible)
    keywordstyle=\color{blue},       % keyword style
    % morekeywords={*,...},
    % if you want to add more keywords to the set
    numbers=left,
    % where to put the line-numbers; possible values are (none, left, right)
    numbersep=5pt,
    % how far the line-numbers are from the code
    numberstyle=\tiny\color{mygray},
    % the style that is used for the line-numbers
    rulecolor=\color{black},
    % if not set, the frame-color may be changed on line-breaks
    % within not-black text (e.g. comments (green here))
    stepnumber=1,
    % the step between two line-numbers.
    % If it's 1, each line will be numbered
    stringstyle=\color{mymauve},     % string literal style
    tabsize=4,                       % sets default tabsize to 4 spaces
    % show the filename of files included with \lstinputlisting;
    % also try caption instead of title
    language = C,
	showspaces = false,
	showtabs = false,
	showstringspaces = false,
	escapechar = ,
}

\def\ContinueLineNumber{\lstset{firstnumber=last}}
\def\StartLineAt#1{\lstset{firstnumber=#1}}
\let\numberLineAt\StartLineAt



\newcommand{\codeline}[1]{
	\alert{\texttt{#1}}
}

% This Document contains the information about this course.

% Authors of the slides
\author{
    Lecturers:
	%name1
	Mirko Jantschke,
	%name2
	Pascal Scholz
	\\
	%creaters of the slides
	Created by: Richard Mörbitz, Manuel Thieme
}


% Fancy Logo
\titlegraphic{\hfill\includegraphics[height=1.25cm]{../templates/fsr_logo_cropped}}


\title{C-Lessons}
\subtitle{Variables}
\date{\today}

\begin{document}

\begin{frame}
	\titlepage
\end{frame}
\begin{frame}{Overview}
	\setbeamertemplate{section in toc}[sections numbered]
	\tableofcontents
\end{frame}

%----------------------------------------------------------------------------------------------------
\begin{frame}[fragile]{Usage}
	Declaration:
	\begin{lstlisting}[numbers=none,basicstyle=\itshape\footnotesize]
type identifier;
\end{lstlisting}
	Assignment: 
	\begin{lstlisting}[numbers=none,basicstyle=\itshape\footnotesize]
identifier = value;
\end{lstlisting}
	Definition (all at once):
	\begin{lstlisting}[numbers=none,basicstyle=\itshape\footnotesize]
type identifier = value;
\end{lstlisting}
	Example:
	\begin{lstlisting}[numbers=none]
int number;                 /* declaration */
number = 42;                /* assignment */
int another_number = 23;    /* definition */

\end{lstlisting}
\end{frame}

%---------------------------------------------------------------------------------------------------
\begin{frame}[fragile]{Saving lines}
	Multiple declarations
	\begin{lstlisting}[numbers=none]
int number, another_number;
\end{lstlisting}
	Multiple Definitions
	\begin{lstlisting}[numbers=none]
int number = 42, anothernumber = 23;
\end{lstlisting}\ \\\ \\
But be careful:
\begin{columns}[c]
	\column{.4\textwidth}
	\begin{lstlisting}[numbers=none]
int a = 23, b = 23;
\end{lstlisting}
	\column{.12\textwidth}
	\centering
	$\neq$
	\column{.3\textwidth}
	\begin{lstlisting}[numbers=none]
int a, b = 23;
\end{lstlisting}
\end{columns}\ \\ \ \\
$\rightarrow$ Avoid multiple variable definitions at one line!
\end{frame}

%----------------------------------------------------------------------------------------------------
\begin{frame}{Valid identifiers}
	\begin{itemize}
		\item Consist of English letters (no \textit{ß}, \textit{ä}, \textit{ö}, \textit{ü}), numbers and underscore (\_)
		\item Start with a letter or underscore
		\item Are case sensitive (\textit{number} differs from \textit{Number})
		\item Must not be reserved words:
	\end{itemize}
    {\scriptsize\tt\begin{tabular}{cccccccc}
        auto	& else	& long	& switch & break	& enum	& register	& typedef\\
        case	& extern	& return	& union & char	& float	& short	& unsigned\\
        const	& for	& signed	& void & continue	& goto	& sizeof	& volatile\\
        default	& if	& static	& while & do	& int	& struct &   double
    \end{tabular}}
    \\\bigskip
	\textbf{Style}:\\
	\begin{itemize}
		\item Stay in one language (English recommended)
		\item Decide whether to use \textit{camelCaseIdentifiers} or \textit{snake\_case\_identifiers}.\\
	\end{itemize}
\end{frame}

%----------------------------------------------------------------------------------------------------
\begin{frame}[fragile]{Speaking identifiers}
	\begin{lstlisting}
/* calculate volume of square pyramid */
int a, b, c;
a = 3;
b = 2;
c = (1 / 3) * a * a * b;
\end{lstlisting}
\centering
$\Downarrow$
	\begin{lstlisting}
/* calculate volume of square pyramid */
int length, height, volume;
length = 3;
height = 2;
volume = (1 / 3) * length * length * height;
\end{lstlisting}
\end{frame}

%----------------------------------------------------------------------------------------------------
\begin{frame}{Use speaking identifiers.}
	\LARGE
	\centering
	Please, use speaking identifiers.\footnotemark
	
	\footnotetext[1]{Seriously, use speaking identifiers.}
\end{frame}

%----------------------------------------------------------------------------------------------------
\section{Data types}
\subsection{}
\begin{frame}[fragile]{Integer numbers}
	\begin{itemize}
		\item Keywords: \textit{int}, \textit{short}, \textit{long}
		\item Stored as a binary number with fixed length
		\item Can be \textit{signed}(default) or \textit{unsigned}
		\item Actual size of \textit{int}, \textit{short}, \textit{long} depends on architecture
	\end{itemize}\ \\
	\ \\
	Example (64 Bit):
	\begin{lstlisting}[numbers=none]
int a;          /* Range: -2.147.483.648 to 2.147.483.647 */
unsigned short b; /* Range: 0 to 65.535 */
\end{lstlisting}
\end{frame}

%---------------------------------------------------------------------------------------------------
\begin{frame}[fragile]{Floating point numbers}
	\begin{itemize}
		\item Keywords: \textit{float}, \textit{double}, \textit{long double}
		\item Stored as specified in \textit{IEEE 754 Standard} TL;DR
		\item Special values for $\infty$, $-\infty$, NaN
		\item Useful for fractions and very large numbers
		\item Type a decimal point instead of a comma!
	\end{itemize}\ \\
	\ \\
	Example:
	\begin{lstlisting}[numbers=none]
float x = 0.125;          /* Precision: 7 to 8 digits */
double y = 111111.111111; /* Precision: 15 to 16 digits */
\end{lstlisting}
\end{frame}
\begin{frame}[fragile]{Characters}
	\begin{itemize}
		\item Keyword: \textit{char}
		\item Can be \textit{signed}(default) or \textit{unsigned}	
		\item Size: 1 Byte (8 Bit) on almost every architecture
		\item Intended to represent a single character
		\item Stores its \textit{ASCII} number (e.g. 'A' $\Rightarrow$ 65)
	\end{itemize}\ \\
	\ \\
	You can define a \textit{char} either by its ASCII number or by its symbol:
	\begin{lstlisting}[numbers=none]
char a = 65;
char b = 'A';       /* use single quotation marks */
\end{lstlisting}
\end{frame}

%----------------------------------------------------------------------------------------------------
\section{Variable I/O}
\subsection{}
\begin{frame}[fragile]{\textit{printf()} with placeholders}
	The string you pass to \textit{printf} may contain placeholders:
	\begin{lstlisting}[numbers=none]
int a = 3;
int b = 5;
float c = 7.4;
printf("a: %d\n", a);
printf("b: %d\nc: %f\n", b, c);
\end{lstlisting}
Output:\begin{lstlisting}[numbers=none]
a: 3
b: 5
c: 7.4
\end{lstlisting}
You can insert any amount of placeholders. For each placeholder, you have to pass a value of the corresponding type.
\end{frame}

%----------------------------------------------------------------------------------------------------
\begin{frame}{Example placeholders}
	The placeholder determines how the value is interpreted.
	To avoid compiler warnings, only use the following combinations: \\ \ \\
	\begin{tabular}{|c|c|c|}
		\hline
		\textbf{type} & \textbf{description} & \textbf{type of argument} \\\hline
		\%c & single character & char, int (if $<=$ 255) \\\hline
		\%d & decimal number & char, int \\\hline
		\%u & unsigned decimal number & unsigned char, unsigned int \\\hline
		\%X & hexadecimal number & char, int \\\hline
		\%ld & long decimal number & long \\\hline
		\%f & floating point number & float, double \\\hline
	\end{tabular}
	\\
	Have a look at the \href{https://github.com/scholzp/c-lessons/blob/master/cheatsheet/cheatsheet.pdf}{Cheatsheet}!
\end{frame}

%----------------------------------------------------------------------------------------------------
\begin{frame}[fragile]{Variable input}
	\textit{scanf()} is another useful function from the standard library.
	\begin{itemize}
		\item Like \textit{printf()}, it is declared in \textit{stdio.h}
		\item Like \textit{printf()}, it has a format string with placeholders
		\item You can use it to read values of primitive datatypes from the command line
	\end{itemize}
	\ \\ \ \\ Example:
	\begin{lstlisting}[numbers=none]
int i;
scanf("%d", &i);
\end{lstlisting}
	After calling \textit{scanf()}, the program waits for the user to input a value in the command line.
	After pressing the \textit{return} key, that value is stored in \textit{i}.
\end{frame}

%----------------------------------------------------------------------------------------------------
\begin{frame}{Note:}
	\begin{itemize}
		\item \textit{scanf()} uses the same placeholders as \textit{printf()}
		\item You must type an \textit{\&} before each variable identifier \\
			(more about this later)
		\item If you read a number (using \%d, \%u etc.), interpretation
		\begin{itemize}
			\item Starts at first digit
			\item Ends before last non digit character
		\end{itemize}
		\item If you use \%c, the first character of the user input is interpreted (this may be a ' ' as well!)
	\end{itemize}
	Never trust the user: they may enter a blank line while you expect a number, which means your input variable is still undefined!
\end{frame}
%------------------------------------------- end ---------------------------------------------------------
\section{Related Task}
\begin{frame}{Task: Rot13}
	\href{http://fsr.github.io/c-lessons/exercises/02_rot_13.html}{Task as online:} \newline
    Rot13 is an ancient cryptography procedure. 
    It encrypts a message by shifting each character 13 letters in alphabetical order (modulo z).\newline

    Write a program that takes a printable character and prints its Rot13 encrypted version. 
    Assume that the input is always a small character. \newline

    Experts: Get it working with keeping the character case (a -\textgreater, B -\textgreater O). 
    Keep it simple. Do not use stuff that has not been introduced yet.

\end{frame}

\end{document}
